\section{Related Work}

\paragraph{Vision-Language-Action Models}

\paragraph{World Model for Embodied Control}

\paragraph{Reinforcement Learning for Vision-Language-Action Models}
Recent work has begun to move beyond pure imitation learning and study reinforcement learning as a post-training mechanism for Vision-Language-Action (VLA) models. iRe-VLA shows that directly applying online RL to large VLA policies is often unstable and computationally demanding, and therefore alternates between RL-based exploration and supervised learning to absorb successful trajectories more stably. RIPT-VLA instead advocates a simpler critic-free interactive post-training paradigm that uses only sparse binary success rewards, combining dynamic rollout sampling with leave-one-out advantage estimation for stable policy updates. To improve long-horizon credit assignment, TGRPO extends group-relative optimization with both step-level and trajectory-level relative advantages, relying on dense multi-stage rewards to provide finer-grained supervision. From a systems perspective, VLA-RL formulates auto-regressive VLA control as trajectory-level online RL and introduces a robotic process reward model, together with implementation strategies such as critic warmup and scalable batch decoding, to make large-scale online optimization practical. Complementary to direct environment interaction, VLA-RFT leverages a data-driven world model as a controllable simulator and fine-tunes VLA policies using verified dense trajectory-level rewards computed from imagined visual rollouts. Overall, these studies suggest three emerging directions for RL-enhanced VLA training: alternating RL with imitation-style updates, direct online post-training under sparse or shaped rewards, and world-model-based reinforcement fine-tuning in learned simulators. 